% Part: incompleteness
% Chapter: arithmetization-syntax
% Section: proofs-in-ax

\documentclass[../../../include/open-logic-section]{subfiles}

\begin{document}

\olfileid{inc}{art}{pax}
\olsection{స్వీకృతాధారిత \usetoken{P}{derivation}}

\begin{explain}
స్వీకృతాధారిత \tetoken{వ్యుత్పత్తులను}{!!{derivation}s}
అంకగణితీకరించాలంటే, \tetoken{వ్యుత్పత్తులను}{!!{derivation}s} సంఖ్యలుగా
ప్రాతినిధ్యం చేయాలి.
\tetoken{వ్యుత్పత్తులు}{!!{derivation}s} కేవలం
\tetoken{సూత్రాల}{!!{formula}s} క్రమాలే కనుక, ప్రతి
\tetoken{వ్యుత్పత్తినీ}{!!{derivation}} అందులోని
\tetoken{సూత్రాల}{!!{formula}s} సంకేతసంఖ్యల క్రమపు సంకేతసంఖ్యగా
సంకేతీకరించడమే సహజమైన పద్ధతి.
\end{explain}

\begin{defn}
$\delta$ అనేది $!A_1$, \dots, $!A_n$ అనే
\tetoken{సూత్రాలతో}{!!{formula}s} కూడిన స్వీకృతాధారిత
\tetoken{వ్యుత్పత్తి}{!!{derivation}} అయితే, $\Gn{\delta}$ అనేది
\[
\tuple{\Gn{!A_1},\dots,\Gn{!A_n}}.
\]
\end{defn}

\begin{ex}
  ఈ అతి సరళమైన \tetoken{వ్యుత్పత్తిని}{!!{derivation}} పరిగణించండి:
  \begin{derivation}
  1. & $!B \lif (!B \lor !A)$ \\
  2. & $(!B \lif (!B \lor !A)) \lif (!A \lif (!B \lif (!B \lor !A)))$\\
  3. & $!A \lif (!B \lif (!B \lor !A))$
  \end{derivation}
  ఈ \tetoken{వ్యుత్పత్తి}{!!{derivation}} గ్యోడెల్ సంఖ్య
  \begin{align*}
  \openTuple\,
    &\Gn{!B\lif(!B\lor !A)},\\
    &\Gn{(!B\lif(!B\lor !A))\lif(!A\lif(!B\lif(!B\lor !A)))},\\
    &\Gn{!A\lif(!B\lif(!B\lor !A))}\,\closeTuple.
  \end{align*}
\end{ex}

\begin{explain}
\tetoken{వ్యుత్పత్తుల}{!!{derivation}s} ప్రాతినిధ్యాన్ని నిర్ణయించిన తరువాత,
అటువంటి \tetoken{వ్యుత్పత్తులను}{!!{derivation}s} ఆదిమ పునరావృత్తంగా
మార్చగలమనీ, వాటి ముఖ్య ధర్మాలు, సంబంధాలను అలాగే వ్యక్తీకరించగలమనీ చూపాలి.
కొన్ని క్రియలు సరళమైనవి: ఉదాహరణకు,
\tetoken{ఒక వ్యుత్పత్తి}{!!a{derivation}} గ్యోడెల్ సంఖ్య~$d$ ఇస్తే,
$(d)_{\len{d}-1}$ దాని అంత్య-\tetoken{సూత్రం}{!!{formula}} గ్యోడెల్
సంఖ్యను ఇస్తుంది. కొన్ని క్రియలు మరింత కష్టం. వాటిని ఎలా చేయాలో కనీసం
రూపురేఖలుగా చూద్దాం. ``$\delta$ అనేది $\Gamma$ నుంచి $!A$ యొక్క
\tetoken{ఒక వ్యుత్పత్తి}{!!a{derivation}}'' అనే సంబంధం $\delta$, $!A$ల
గ్యోడెల్ సంఖ్యలలో ఆదిమ పునరావృత్తమని చూపడమే లక్ష్యం.
\end{explain}

\begin{prop}
\ollabel{prop:followsby}
క్రింది సంబంధాలు ఆదిమ పునరావృత్తం:
\begin{enumerate}
\item $!A$ ఒక స్వీకృత రూపం.
\item $\delta$లో $i$-వ పంక్తి మోడస్ పోనెన్స్ ద్వారా సమర్థించబడింది.
\item $\delta$లో $i$-వ పంక్తి \QR{} ద్వారా సమర్థించబడింది.
\item $\delta$ సరైన \tetoken{వ్యుత్పత్తి}{!!{derivation}}.
\end{enumerate}
\end{prop}

\begin{proof}
\tetoken{సూత్రాల}{!!{formula}s} గ్యోడెల్ సంఖ్యలకూ,
\tetoken{వ్యుత్పత్తుల}{!!{derivation}s} గ్యోడెల్ సంఖ్యలకూ మధ్య అనురూపమైన
సంబంధాలు ఆదిమ పునరావృత్తమని చూపాలి.
\begin{enumerate}
\item మనకు స్వీకృత రూప పథకాల నిర్ణీత జాబితా ఉంది; $!A$ వాటిలో ఒక పథకం
  ఇచ్చే రూపంలో ఉంటే అది స్వీకృత రూపం. పథకాల జాబితా పరిమితం కనుక, ప్రతి
  స్వీకృత రూప పథకానికీ $!A$ ఆ రూపంలో ఉందా అని ఆదిమ పునరావృత్తంగా
  పరీక్షించగలమని చూపితే చాలు. ఉదాహరణకు ఈ స్వీకృత రూప పథకాన్ని పరిగణించండి:
  \[
  !B\lif(!C\lif !B).
  \]
  \tetoken{సూత్రాలు}{!!{formula}s} $!B$, $!C$ ఉండి, `$($'కు $!B$ను,
  దానికి `$\lif$'ను, దానికి `$($'ను, దానికి $!C$ను, దానికి `$\lif$'ను,
  దానికి $!B$ను, చివరికి `$))$'ను సంయోజించినప్పుడు $!A$ వస్తే, $!A$ ఈ
  స్వీకృత రూప పథకానికి ఒక నిదర్శనం. $!A$ గ్యోడెల్ సంఖ్య~$n$కు అనురూపమైన
  ధర్మాన్ని పరీక్షించవచ్చు: క్రమాల సంయోజనం ఆదిమ పునరావృత్తం; ఈ సంబంధం
  నిలిచినప్పుడు $!B$, $!C$ రెండూ $!A$కు ఉప-\tetoken{సూత్రాలు}{!!{formula}s}.
  కాబట్టి $!B$, $!C$ల గ్యోడెల్ సంఖ్యలు $!A$ గ్యోడెల్ సంఖ్యకంటే తక్కువగా
  ఉండాలి.
  అందువల్ల ఇలా నిర్వచించవచ్చు:
  \begin{multline*}
  \fn{IsAx}_{!B\lif(!C\lif !B)}(n)\defiff\bexists{b<n}{
    \bexists{c<n}{(\fn{Sent}(b)\land\fn{Sent}(c)\land{}}}\\ n=
  \Gn{(}\concat b\concat\Gn{\lif}\concat\Gn{(}\concat c\concat
  \Gn{\lif}\concat b\concat\Gn{))}).
  \end{multline*}
  ప్రతి స్వీకృత రూప పథకానికీ అటువంటి నిర్వచనం ఉంటే, వాటి వికల్పం
  ``$n$ ఒక స్వీకృత రూపపు గ్యోడెల్ సంఖ్య'' అనే $\fn{IsAx}(n)$ ధర్మాన్ని
  నిర్వచిస్తుంది.
\item $\delta$లో $i$-వ పంక్తి మోడస్ పోనెన్స్ ద్వారా సమర్థించబడటానికి
  అవసరమైనదీ, సరిపడేదీ: $j$ అనే పంక్తి, $k<i$ అనే పంక్తి ఉండి, $j$-వ పంక్తిలోని
  \tetoken{వాక్యం}{!!{sentence}} ఏదో సూత్రం~$!A$, $k$-వ పంక్తిలోని
  వాక్యం $!A\lif !B$, $i$-వ పంక్తిలోని వాక్యం $!B$ కావడం.
  \begin{multline*}
    \fn{MP}(d,i)\defiff\bexists{j<i}{\bexists{k<i}{}}\\
    (d)_k=\Gn{(}\concat(d)_j
      \concat\Gn{\lif}\concat(d)_i\concat\Gn{)}
  \end{multline*}
  పరిమిత పరిమాణీకరణ, సంయోజనం, $=$ అన్నీ ఆదిమ పునరావృత్తం కనుక, ఇది
  ఒక ఆదిమ పునరావృత్త సంబంధాన్ని నిర్వచిస్తుంది.
\item $\delta$లో ఒక పంక్తి \QR{} ద్వారా సమర్థించబడటానికి అవసరమైనదీ,
  సరిపడేదీ: అది $!B\lif\lforall[x][!A(x)]$ రూపంలో ఉండటం, ముందరి పంక్తి
  ఏదో \tetoken{స్థిరసంకేతం}{!!{constant}}~$c$కు $!B\lif !A(c)$ కావడం,
  $c$ అనేది $!B$లో కనిపించకపోవడం. అంటే, క్రింది షరతులు నెరవేరాలి:
  \begin{enumerate}
  \item \tetoken{ఒక వాక్యం}{!!a{sentence}}~$!B$ ఉంది;
  \item సరిగ్గా ఒక్క చరం~$x$ స్వేచ్ఛగా గల
    \tetoken{ఒక సూత్రం}{!!a{formula}}~$!A(x)$ ఉంది;
  \item $i$-వ పంక్తిలో $!B\lif\lforall[x][!A(x)]$ ఉంది;
  \item ఏదో స్థిరసంకేతం~$c$కు, ఏదో $j<i$ పంక్తిలో
    $!B\lif\Subst{!A}{c}{x}$ ఉంది;
  \item అది $!B$లో కనిపించదు.
  \end{enumerate}
  వీటన్నింటినీ ఆదిమ పునరావృత్తంగా పరీక్షించవచ్చు. $!B$, $!A(x)$, $x$ల
  గ్యోడెల్ సంఖ్యలు $i$-వ పంక్తిలోని సూత్రపు గ్యోడెల్ సంఖ్యకంటే తక్కువ;
  $c$ గ్యోడెల్ సంఖ్య $j$-వ పంక్తిలోని సూత్రపు గ్యోడెల్ సంఖ్యకంటే తక్కువ:
  \begin{multline*}
    \fn{QR}_1(d,i)\defiff\bexists{j<i}{\bexists{b<(d)_i}{\bexists{x<(d)_i}{
        \bexists{a<(d)_i}{\bexists{c<(d)_j}{(}}}}}\\
    \fn{Var}(x)\land\fn{Const}(c)\land{}\\
    (d)_i=\Gn{(}\concat
    b\concat\Gn{\lif}\concat\Gn{\lforall}\concat x\concat a
    \concat\Gn{)}\land{}\\
    (d)_j=\Gn{(}\concat b\concat
    \Gn{\lif}\concat\fn{Subst}(a,c,x)\concat\Gn{)}\land{}\\
    \fn{Sent}(b)\land\fn{Sent}(\fn{Subst}(a,c,x))\land{}
    \bforall{k<\len{b}}{(b)_k\neq(c)_0})
  \end{multline*}
  \sourcecorrection{OLTEART-012}{ఆంగ్ల మూలంలోని $\fn{QR}_1(d,i)$ సూత్రం $(d)_j$ను ఉపయోగించినా $j$ను పరిమాణీకరించదు. ముందరి పంక్తి కావాలనే వివరణకు అనుగుణంగా $\bexists{j<i}{}$ను చేర్చాం. అదే వివరణలో స్థిరసంకేతం $c$ బదులు ``$a$ యొక్క'' గ్యోడెల్ సంఖ్య $j$-వ పంక్తికంటే తక్కువని వచ్చిన అక్షరదోషాన్ని కూడా $c$గా సరిచేశాం.}
  ఇక్కడ $c$, $x$లు వరుసగా పదాలుగా పరిగణించిన స్థిరసంకేతం, చరం యొక్క
  గ్యోడెల్ సంఖ్యలు (వాటి సంకేతసంఖ్యలు కావు) అని అనుకుంటాం.
  $\Subst{!A(x)}{c}{x}$ \tetoken{ఒక వాక్యం}{!!a{sentence}} అవుతుందా అని
  పరీక్షించి, $x$ అనేది $!A(x)$లోని ఏకైక స్వేచ్ఛా చరమని పరీక్షిస్తాం;
  $!B$లోని ప్రతి సంకేతం $c$కు భిన్నమని కోరడం ద్వారా $c$ అనేది $!B$లో
  కనిపించదని నిర్ధారిస్తాం.

  \QR{} యొక్క మరో రూపాన్ని అభ్యాసంగా వదిలేస్తాం.
\item $d$ ఒక సరైన \tetoken{వ్యుత్పత్తి}{!!{derivation}} గ్యోడెల్ సంఖ్య
  కావడానికి అవసరమైనదీ, సరిపడేదీ: అందులోని ప్రతి పంక్తి స్వీకృత రూపం కావడం,
  లేదా మోడస్ పోనెన్స్ లేదా \QR{} ద్వారా సమర్థించబడటం. కాబట్టి
  \[
  \fn{Deriv}(d)\defiff\bforall{i<\len{d}}{(\fn{IsAx}((d)_i)\lor
    \fn{MP}(d,i)\lor\fn{QR}(d,i))}
  \]
\end{enumerate}
\end{proof}

\begin{prob}
\olref[inc][art][pax]{prop:followsby}లోలాగే ఈ సంబంధాలను నిర్వచించండి:
\begin{enumerate}
\item $\fn{IsAx}_{!A\lif(!B\lif(!A\land !B))}(n)$,
\item $\fn{IsAx}_{\lforall[x][!A(x)]\lif !A(t)}(n)$,
\item $\fn{QR}_{2}(d,i)$ (\QR{} యొక్క మరో రూపానికి).
\end{enumerate}
\end{prob}

\begin{prop}
$\Gamma$ అనేది \tetoken{వాక్యాల}{!!{sentence}s} ఆదిమ పునరావృత్త సమితి
అనుకుందాం. ``$x$ అనేది $\Gamma$ నుంచి $!A$ యొక్క
\tetoken{ఒక వ్యుత్పత్తి}{!!a{derivation}}~$\delta$ సంకేతసంఖ్య, $y$ అనేది
$!A$ గ్యోడెల్ సంఖ్య'' అని వ్యక్తీకరించే $\Prf[\Gamma](x,y)$ సంబంధం
ఆదిమ పునరావృత్తం.
\end{prop}

\begin{proof}
``$y\in\Gamma$''ను ఆదిమ పునరావృత్త విధేయం $R_\Gamma(y)$ ఇస్తుందని
అనుకుందాం. $y$ అనేది \tetoken{ఒక వాక్యం}{!!a{sentence}}~$!A$ గ్యోడెల్
సంఖ్యగా, $x$ అనేది $\Gamma$ నుంచి $!A$ యొక్క
\tetoken{ఒక వ్యుత్పత్తి}{!!a{derivation}} సంకేతసంఖ్యగా ఉన్నప్పుడే నిలిచే
$\Prf[\Gamma](x,y)$ సంబంధం ఆదిమ పునరావృత్తమని, అంటే
$\Prf[\Gamma](x,y)$ నిలిచే షరతు ఆదిమ పునరావృత్తమని చూపాలి.

మునుపటి ప్రతిపాదన ప్రకారం, $x$ సరైన
\tetoken{వ్యుత్పత్తి}{!!{derivation}}~$\delta$ సంకేతసంఖ్య అయితేనే నిలిచే
$\fn{Deriv}(x)$ ధర్మం ఆదిమ పునరావృత్తం. అయితే ఆ నిర్వచనం
\tetoken{వ్యుత్పత్తిలో}{!!{derivation}} పంక్తిని
సమర్థించే అదనపు మార్గంగా $\Gamma$ సమితిని పరిగణించలేదు. ఒక పంక్తి
\QR{} ద్వారా సమర్థించబడిందా అనే మన ఆదిమ పునరావృత్త పరీక్షలో, స్థిరసంకేతం~$c$
$\Gamma$లో కనిపించకూడదనే ఆవశ్యకతను కూడా పరిగణించలేదు. $\Gamma$ను నేరుగా
పరిగణించేలా నిర్వచనాన్ని సవరించవచ్చు; కానీ $\fn{Deriv}$ను, నిగమన సిద్ధాంతాన్ని
వాడటం సులభం. $\Gamma\Proves !A$ కావడానికి అవసరమైనదీ, సరిపడేదీ:
$!B_1$, \dots, $!B_n\in\Gamma$ అనే ఏదో పరిమిత
\tetoken{వాక్యాల}{!!{sentence}s} జాబితా ఉండి,
$\{!B_1,\dots,!B_n\}\Proves !A$ కావడం. నిగమన సిద్ధాంతం ప్రకారం,
$\Proves(!B_1\lif(!B_2\lif\cdots(!B_n\lif !A)\cdots))$ అయినప్పుడే ఇది
జరుగుతుంది. గ్యోడెల్ సంఖ్య~$z$ గల \tetoken{ఒక వాక్యం}{!!a{sentence}}
ఈ రూపంలో ఉందా అని ఆదిమ పునరావృత్తంగా పరీక్షించవచ్చు. కాబట్టి $x$ను
$\Gamma$ నుంచి గ్యోడెల్ సంఖ్య~$y$ గల \tetoken{వాక్యం}{!!{sentence}}
యొక్క \tetoken{ఒక వ్యుత్పత్తి}{!!a{derivation}} గ్యోడెల్ సంఖ్యగా తీసుకునే
బదులు, $x$ను పై రూపంలోని అంతర్నిహిత షరతువాక్యానికి $\emptyset$ నుంచి వచ్చిన
\tetoken{ఒక వ్యుత్పత్తి}{!!a{derivation}} గ్యోడెల్ సంఖ్యగా తీసుకుంటాం.

మొదట, \tetoken{వాక్యాల}{!!{sentence}s} క్రమం ఇస్తే, ఆ వాక్యాలన్నీ
పూర్వపక్షాలుగా, ఇచ్చిన \tetoken{వాక్యం}{!!{sentence}} పర్యవసానంగా గల
షరతువాక్యాన్ని ఆదిమ పునరావృత్తంగా నిర్మించవచ్చు:
\begin{align*}
  \fn{hCond}(s,y,0)&=y\\
  \fn{hCond}(s,y,n+1)&=\Gn{(}\concat(s)_n\concat\Gn{\lif}
  \concat\fn{hCond}(s,y,n)\concat\Gn{)}\\
  \fn{Cond}(s,y)&=\fn{hCond}(s,y,\len{s})\\
  \intertext{కాబట్టి $\Prf[\Gamma](x,y)$ను ఇలా నిర్వచించవచ్చు:}
  \Prf[\Gamma](x,y)&\defiff\bexists{s<\fn{sequenceBound}(x,x)}{(}\\
  &\qquad(x)_{\len{x}-1}=\fn{Cond}(s,y)\land{}\\
  &\qquad\bforall{i<\len{s}}{(s)_i\in\Gamma}\land{}\\
  &\qquad\fn{Deriv}(x)).
\end{align*}
\sourcecorrection{OLTEART-013}{ఆంగ్ల మూలంలోని ఆదిమ పునరావృత్త దశ $\fn{hCond}(s,y,n+1)$ను నిర్వచిస్తూ, మూడు ఆర్గుమెంట్లతో నిర్వచించని $\fn{Cond}(s,y,n)$ను పిలుస్తుంది. ఉద్దేశించిన ముందరి పునరావృత్త విలువ $\fn{hCond}(s,y,n)$ను వాడాం.}
$s$పై పరిమితి ఇలా లభిస్తుంది: ప్రతి $(s)_i$ అనేది
\tetoken{వ్యుత్పత్తి}{!!{derivation}} చివరి పంక్తికి ఉప-
\tetoken{సూత్రం}{!!{formula}} యొక్క గ్యోడెల్ సంఖ్య; కాబట్టి అది
$(x)_{\len{x}-1}$కంటే తక్కువ. పూర్వపక్షాలు $!B\in\Gamma$ల సంఖ్య,
అంటే $s$ పొడవు, $x$ చివరి పంక్తి పొడవుకంటే తక్కువ.
\end{proof}

\end{document}
